\documentclass{article}
\usepackage{mystyle}

\usetikzlibrary{arrows.meta, positioning, calc, fit, backgrounds}

\definecolor{ink}{HTML}{1F2933}
\definecolor{muted}{HTML}{667085}
\definecolor{blue}{HTML}{2F80ED}
\definecolor{line}{HTML}{D0D5DD}
\definecolor{soft}{HTML}{F7F9FC}
\definecolor{cell}{HTML}{FFFFFF}
\definecolor{mark}{HTML}{EBF3FF}
\definecolor{green}{HTML}{12B76A}
\definecolor{orange}{HTML}{F79009}
\definecolor{purple}{HTML}{7A5AF8}
\definecolor{redsoft}{HTML}{F97066}

\tikzset{
  paperbox/.style={
    draw=ink!55,
    fill=white,
    rounded corners=2pt,
    very thick,
    minimum height=10mm,
    text width=25mm,
    align=center,
    font=\small
  },
  papersmall/.style={
    draw=ink!45,
    fill=white,
    rounded corners=2pt,
    thick,
    minimum height=8mm,
    text width=22mm,
    align=center,
    font=\small
  },
  papernote/.style={
    font=\footnotesize,
    text=muted,
    align=center
  },
  tensorlabel/.style={
    font=\footnotesize\ttfamily,
    text=ink
  },
  paperarrow/.style={
    -{Latex[length=2.2mm]},
    thick,
    draw=ink!75
  },
  looparrow/.style={
    -{Latex[length=2.2mm]},
    thick,
    draw=ink!75
  }
}

\newcommand{\filtergrid}[4]{%
\begin{tikzpicture}[x=8mm,y=8mm,baseline=-10mm]
  \foreach \i in {0,...,3} {
    \draw[line] (0,-\i) -- (3,-\i);
    \draw[line] (\i,0) -- (\i,-3);
  }

  \fill[fill=mark] (#2,-#3) rectangle ++(1,-1);
  \fill[fill=mark] (1,-1) rectangle ++(1,-1);

  \node at (#2+.5,-#3-.5) {$1$};
  \node at (1.5,-1.5) {$-2$};

  \node[font=\small, text=ink] at (1.5,.45) {$#1$};
  \node[font=\scriptsize, text=muted] at (1.5,-3.45) {#4};
\end{tikzpicture}%
}

\DeclareMathOperator{\ReLU}{ReLU}

\title{\vspace{-1.2em}5주차: 시각화 고고학}
\author{ML기초 시그 by 윤승현}
\date{}

\begin{document}
\maketitle

\subsection*{정글 속 디스크}

정글에 불시착한 yogurt는 울창한 밀림에서 헤매고 있다.
끝없이 이어진 밀림 사이의 강길을 따라 이동하던 중,
먼지가 잔뜩 묻은 디스크 하나를 발견한다.

\begin{quote}
\centering
\itshape
이 퍼즐을 푼 자, 자신이 이 숲을 살아나갈 수 있는지 알 수 있을 것이다.
\end{quote}

식량은 고작 10일 치.
앞길은 막막하다.
그런데도 yogurt는 이상하게 이 문제에 끌린다.

자세히 퍼즐을 읽어보니 다음과 같이 적혀 있었다.

\begin{center}
\fbox{\parbox{0.82\linewidth}{
\centering
$\{0,1,2\}$로 채워진 $11 \times 11$ 이미지를 입력으로 받는다.\\
이 신경망이 무슨 역할을 수행하는지 알아맞혀라.
}}
\end{center}

\subsection*{단서 1: 필터}

디스크에는 다음과 같은 네 개의 $3 \times 3$ 필터가 함께 들어 있었다.

\begin{center}
\begin{tabular}{@{}c@{\qquad}c@{\qquad}c@{\qquad}c@{}}
\filtergrid{K_{\leftarrow}}{0}{1}{왼쪽}
&
\filtergrid{K_{\uparrow}}{1}{0}{위쪽}
&
\filtergrid{K_{\rightarrow}}{2}{1}{오른쪽}
&
\filtergrid{K_{\downarrow}}{1}{2}{아래쪽}
\end{tabular}
\end{center}

\subsection*{단서 2: 네트워크 구조}

신경망의 전체 구조는 다음과 같다.

\vspace{4mm}

\resizebox{\linewidth}{!}{
\begin{tikzpicture}[
  >=Latex,
  font=\small,
  node distance=8mm,
  arr/.style={
    -{Latex[length=2.2mm,width=1.6mm]},
    line width=0.55pt,
    draw=black!75
  },
  darr/.style={
    -{Latex[length=2.2mm,width=1.6mm]},
    dashed,
    line width=0.55pt,
    draw=black!45
  },
  skiparr/.style={
    -{Latex[length=2.0mm,width=1.5mm]},
    line width=0.55pt,
    draw=black!65
  },
  box/.style={
    rounded corners=3.5pt,
    line width=0.55pt,
    draw=black!45,
    fill=white,
    align=center,
    inner sep=3pt
  },
  topbox/.style={
    box,
    minimum width=34mm,
    minimum height=13mm
  },
  mainbox/.style={
    box,
    draw=blue!55!black,
    fill=blue!4,
    minimum width=45mm,
    minimum height=14mm,
    line width=0.8pt
  },
  module/.style={
    box,
    minimum height=10mm,
    minimum width=20mm
  },
  tensor/.style={
    module,
    draw=green!45!black,
    fill=green!10,
    minimum width=16mm
  },
  conv/.style={
    module,
    draw=redsoft!70!black,
    fill=redsoft!10,
    minimum width=27mm
  },
  act/.style={
    module,
    draw=orange!60!black,
    fill=orange!10,
    minimum width=18mm
  },
  gate/.style={
    module,
    draw=purple!60!black,
    fill=purple!8,
    minimum width=18mm
  },
  reduce/.style={
    module,
    draw=violet!55!black,
    fill=violet!9,
    minimum width=29mm
  },
  clip/.style={
    module,
    draw=orange!60!black,
    fill=orange!11,
    minimum width=18mm
  },
  plus/.style={
    circle,
    draw=black!70,
    fill=white,
    line width=0.55pt,
    minimum size=6.8mm,
    inner sep=0pt,
    font=\normalsize
  },
  label/.style={
    font=\scriptsize,
    text=black!55,
    align=center
  },
  paneltitle/.style={
    font=\bfseries\small,
    text=blue!55!black
  }
]

% =====================================================
% Top-level unrolled architecture
% =====================================================

\node[topbox] (input) {
  \textbf{Input image}\\[-0.5mm]
  $X_0 \in \{0,1,2\}^{11\times 11}$
};

\node[mainbox, right=18mm of input] (repeat) {
  \textbf{ConvBlock} $\times 10$
};

\node[topbox, right=18mm of repeat, minimum width=34mm] (out) {
  $(X_{10})_{11,11}=2$
};

\draw[arr] (input) -- (repeat);
\draw[arr] (repeat) -- (out);

% =====================================================
% Expanded ConvBlock
% =====================================================

\node[tensor, below=35mm of repeat, xshift=-68mm] (xt) {$X_t$};

\node[conv, right=6mm of xt] (convop) {
  $3{\times}3$ Conv\\[-0.4mm]
  $4$ kernels
};

\node[act, right=6mm of convop] (relu) {$\ReLU$};

\node[gate, right=6mm of relu] (rho) {$\rho_2$};

\node[reduce, right=6mm of rho] (sum) {
  channel sum\\[-0.4mm]
  $[4,H,W]\to[H,W]$
};

\node[plus, right=6mm of sum] (add) {$+$};

\node[clip, right=6mm of add] (sat) {$\text{clip}_{[0,2]}$};

\node[tensor, right=6mm of sat] (xtp) {$X_{t+1}$};

\draw[arr] (xt) -- (convop);
\draw[arr] (convop) -- (relu);
\draw[arr] (relu) -- (rho);
\draw[arr] (rho) -- (sum);
\draw[arr] (sum) -- (add);
\draw[arr] (add) -- (sat);
\draw[arr] (sat) -- (xtp);

% Residual identity path
\coordinate (s1) at ($(xt.north)+(0,8mm)$);
\coordinate (s2) at ($(add.north)+(0,8mm)$);

\draw[skiparr]
  (xt.north) -- (s1)
  -- node[label, above=0.6mm] {identity}
  (s2) -- (add.north);

% Panel background, drawn behind nodes
\begin{pgfonlayer}{background}
\node[
  rounded corners=6pt,
  draw=blue!55!black,
  fill=blue!2,
  line width=0.75pt,
  inner xsep=7mm,
  inner ysep=7mm,
  fit=(xt) (convop) (relu) (rho) (sum) (add) (sat) (xtp) (s1) (s2)
] (panel) {};
\end{pgfonlayer}

\node[paneltitle, anchor=north west]
  at ([xshift=4mm,yshift=-3mm]panel.north west) {ConvBlock};

% Zoom relation
\draw[darr]
  (repeat.south) -- node[label, right=1mm] {one step} (panel.north);

% Tiny annotations
\node[label, below=1mm of xt] {$[H,W]$};
\node[label, below=1mm of sum] {aggregate};
\node[label, below=1mm of add] {residual};

\end{tikzpicture}
}

\begin{center}
\fbox{\parbox{0.84\linewidth}{
\centering
질문: 이 네트워크는 무엇을 하는가?
}}
\end{center}

\end{document}